\section{Discussion}

\paragraph{Temporal structure as the primary discriminator.}
The performance gap between Trajectory Emb (C\,=\,0.624) and Disease
Text Emb (C\,=\,0.616) --- with identical encoders and raw features --- is
the cleanest signal in our results.
Prose serialisation collapses temporal ordering into semantic proximity:
``At age 44, J45'' and ``At age 52, F32'' appear near each other in
embedding space regardless of the 8-year interval.
Token-stream serialisation (``44yr:J45 $\rightarrow$ 52yr:F32'') preserves
ordinal structure that the encoder translates into positional context.
This finding aligns with the structured multimodal intelligence thesis:
surface semantic encoding is insufficient when temporal causal structure
drives the outcome.

\paragraph{LLMs can encode numeric structure from prose.}
The Clinical Summary Emb result is perhaps the most surprising: Qwen3-8B
matches a 39-feature CoxPH using only prose descriptions of the same
features.
This suggests that general-purpose LLM pre-training has internalised
enough biomarker-range semantics to distinguish, e.g., HbA1c\,=\,52 from
HbA1c\,=\,38 in terms of mortality relevance.
The near-term TD-AUC advantage (0.605 vs.\ 0.584 at 9.5\,yr) further
suggests the embedding captures interactions between biomarker values and
disease history that a linear CoxPH on raw features misses.
However, this numeric encoding is \emph{implicit} --- it cannot be
inspected or enforced --- which limits trustworthiness in high-stakes
clinical deployment.

\paragraph{The calibration gap of zero-shot inference.}
Delphi's IBS of 0.185 vs.\ $\approx$0.141 for all CoxPH methods reveals
a structural grounding problem: the sigmoid hazard chain produces survival
curves that are systematically miscalibrated relative to observed event
rates.
The near-term AUC collapse (0.52 at 9.5\,yr) suggests the model's
internal hazard estimates do not properly weight acute risk factors
(CRP, creatinine) that dominate near-term mortality.
Fine-tuning with a proper scoring rule (e.g., Brier score loss) would
likely close this gap, but at the cost of zero-shot generalisability.

\paragraph{Limitations.}
Our embedding methods freeze Qwen3-8B; fine-tuning may yield additional
gains at substantially higher compute cost.
IBS and TD-AUC rely on independent censoring assumptions that may be
violated in UKB.
All results are from a single biobank cohort; external validation is
required before clinical translation.
